% ============================================================
% PRE-REVIEW NOTES — REUSABLE TEMPLATE
% ============================================================
% Domain-agnostic. Fill in the placeholder content marked [LIKE THIS]
% for any subject: math, CS, chemistry, law, medicine, history, etc.
%
% Design goals of this template:
%  - One consistent visual language across every chapter of a course,
%    so a learner's full notes set reads as one coherent book.
%  - Every section carries a stable source ID and a precise locator.
%    Use \sourceref for teaching material and \assessref for evidence
%    that a topic is assessed.
%  - A \note{} macro for optional inline glosses in a second language,
%    left empty / unused if not needed.
%  - Structural boxes for: derivations (\rc), memorable takeaways
%    (\keybox), and numbered steps (\step) inside a derivation.
% ============================================================

\documentclass[a4paper,10pt]{article}

\usepackage[utf8]{inputenc}
\usepackage[left=2.2cm, right=1.2cm, top=1.2cm, bottom=1.2cm]{geometry}
\usepackage{xcolor}
\usepackage{amsmath, amssymb}
\usepackage{tikz}
\usepackage{pifont}
\usepackage{tcolorbox}
\usepackage{hyperref}
\usepackage{booktabs}
\usepackage{everypage}

% OPTIONAL: bilingual annotation support.
% Uncomment the next line (and the \cn / \note macros below that use it)
% only if the learner wants inline glosses in a second language.
% Requires a CJK-capable engine (xelatex/lualatex) or the CJK package
% for Chinese/Japanese/Korean specifically.
% \usepackage{CJKutf8}

\tcbuselibrary{skins, breakable}

% ---------- Colour palette (change freely per learner's taste) ----------
\definecolor{barblue}{HTML}{2E75B6}
\definecolor{titleblue}{HTML}{2E75B6}
\definecolor{redline}{HTML}{C0392B}
\definecolor{yellowbg}{HTML}{FFF3CD}
\definecolor{yellowtext}{HTML}{856404}
\definecolor{darktext}{HTML}{2C3E50}
\definecolor{lightblue}{HTML}{D6EAF8}

\pagestyle{empty}
\setlength{\parindent}{0pt}

% ---------- OPTIONAL bilingual gloss macro ----------
% If CJK is enabled above, use: \begin{CJK}{UTF8}{gbsn}#1\end{CJK}
% Otherwise this just renders the text in italics — safe default.
\newcommand{\cn}[1]{\textit{#1}}

% ---------- Structural macros ----------
\newcommand{\chaptertitle}[1]{{\large\bfseries\color{titleblue} #1}\\[2pt]}
% \pref: tag showing WHERE in the source materials this chapter/section
% lives (e.g. "Lec W4", "Ch. 7", "Module 3"). Purely a locator, distinct
% from \examref which tags WHERE it is ASSESSED.
\newcommand{\pref}[1]{\colorbox{yellowbg}{\scriptsize\texttt{\color{yellowtext}#1}}}
\newcommand{\subtitle}[1]{{\scriptsize\color{gray}\itshape #1}}

\newcommand{\topicsection}[1]{%
  \vspace{6pt}{\noindent\normalsize\bfseries\color{titleblue} #1}\\[-6pt]
  {\color{barblue}\rule{\linewidth}{1.2pt}}\\[2pt]
}

% \examref: the assessment-mapping tag. Put worksheet numbers, assignment
% numbers, exam year+question, or "Sample Solutions" here. This is the
% single most important tag in the whole template — it is what lets a
% learner instantly see "this concept is worth my time" and go find the
% matching practice problem.
\reversemarginpar
\setlength{\marginparwidth}{1.8cm}
\setlength{\marginparsep}{2pt}
\newcommand{\examref}[1]{\marginpar{\raggedleft\tiny\bfseries\color{redline}\texttt{#1}}}

% Machine-checkable citation wrappers. Source IDs come from the JSON
% material inventory generated by scripts/inventory_materials.py.
\newcommand{\sourceref}[2]{\pref{#1: #2}}
\newcommand{\assessref}[2]{\examref{#1: #2}}

% \rc: a worked-derivation / detailed-content box with a red left border.
% Use for concept explanations AND for full worked derivations of real
% past problems (the worked-example stage of the workflow).
\newtcolorbox{rc}{
  blanker, breakable,
  borderline west={2pt}{0pt}{redline},
  left=8pt, top=1pt, bottom=1pt,
  before skip=0pt, after skip=3pt,
}

% \keybox: a highlighted takeaway — the one sentence a learner should
% remember from this section even if they forget everything else.
\newtcolorbox{keybox}{
  colback=lightblue, colframe=barblue,
  boxrule=0.8pt, arc=2pt,
  left=4pt, right=4pt, top=2pt, bottom=2pt,
  before skip=4pt, after skip=4pt,
}

% \step{①}{text}: a numbered step inside a derivation. Use circled
% numbers (Pifont \ding{192}-\ding{201} for ①-⑩) or bold digits.
\newcommand{\step}[2]{{\small\color{redline}\bfseries #1}\;{\small #2}\\[1pt]}

% \note{}: optional inline personal gloss (e.g. second-language
% annotation, or a personal mnemonic). Leave unused if not needed.
\newcommand{\note}[1]{{\footnotesize\color{titleblue}\itshape #1}\\[1pt]}

% Coloured bar across the top of every page — purely cosmetic branding
% for the notes set; delete the \AddEverypageHook block if unwanted.
\AddEverypageHook{\tikz[remember picture, overlay]{
  \fill[barblue] ([yshift=0pt]current page.north west) rectangle ([yshift=-14pt]current page.north east);
}}

\begin{document}

% ============================================================
% CHAPTER HEADER
% ============================================================
\vspace*{2pt}
\chaptertitle{Chapter [X]: [Topic Name]}
\sourceref{S001}{[e.g. PDF p. 18 / Section 4]} \quad
\subtitle{[One line: why this chapter matters for assessment — e.g.
"18 of 100 exam marks; recurs in every past sitting; official sample
solution exists for Problem 3." If unsupported, state that the priority
is unknown rather than filling this placeholder with a generic claim.]}

\vspace{4pt}

% ============================================================
% CONCEPT SECTION — repeat this block for each core idea
% ============================================================
\topicsection{[X.1] [Concept Name]}
\sourceref{S001}{[where taught]}\assessref{S004}{[e.g. 2023 exam Q5]}

\begin{rc}
{\small\bfseries [One plain-language framing sentence: what problem
does this concept solve, or what question does it answer?]}

\medskip
{\small\bfseries Formal statement:}\;
[Definition / theorem / rule / doctrine, in the course's own notation
and phrasing — not a generic textbook restatement.]

\medskip
{\small\bfseries Short example:}\;
[A minimal example illustrating the concept — this is NOT the full
worked derivation; that goes in its own "Derivation" block below.]

\begin{keybox}
\centering\small\bfseries [The one sentence to remember from this section.]
\end{keybox}
\end{rc}

% ============================================================
% WORKED DERIVATION — the highest-value block in the whole template.
% Use a REAL problem from the learner's own worksheets/assignments/
% past exams here, not an invented analogue. Include every step a
% full-credit answer would need, calibrated against any official
% sample solution's depth and notation.
% ============================================================
\topicsection{Derivation: [Descriptive title, e.g. "2022 Exam Q4 Walkthrough"]}
\sourceref{S002}{[problem location]}\assessref{S004}{[exam year + Q#, if applicable]}

\begin{rc}
{\small\bfseries Problem statement:}\; [Restate the real problem concisely.]

\medskip
{\small\bfseries Solution:}
\step{\ding{192}}{[First step of the derivation, with reasoning.]}
\step{\ding{193}}{[Second step.]}
\step{\ding{194}}{[Continue until the answer is reached — match the
depth/notation of any official sample solution if one exists.]}

\medskip
{\small\bfseries Why this matters:}\;
[One or two sentences on what makes this problem TYPE recur — e.g.
"This exact reasoning pattern appears in 3 of the last 4 exam sittings."]

\begin{keybox}
\centering\small\bfseries [The reusable technique/pattern this
derivation demonstrates — phrased so it transfers to similar problems.]
\end{keybox}
\end{rc}

% ============================================================
% TRAP BOX — a specific, sourced common mistake. Do not invent
% generic "students often confuse X and Y" content — this must come
% from real instructor emphasis, examiner-report language, or a
% pattern actually observed across past assessments.
% ============================================================
\topicsection{\color{redline}\ding{72} Trap: [Name of the specific
confusion or mistake]}
\sourceref{S003}{[source of this warning, with page/section/question]}

\begin{rc}
{\small\bfseries The trap:}\; [What learners commonly get wrong, and why
it looks right at first.]

\medskip
{\small\bfseries The correct reasoning:}\; [What distinguishes the
correct approach, stated precisely.]

\begin{keybox}
\centering\small\bfseries\color{redline} [A short, memorable
correction — the thing to say to yourself in the exam.]
\end{keybox}
\end{rc}

% ============================================================
% QUICK REFERENCE CARD — one dense screen, every formula/definition/
% technique in the chapter, for last-minute review.
% ============================================================
\topicsection{[X.N] Quick Reference Card}

\begin{rc}
{\small\bfseries [Concept 1]:}\; [one-line formula/definition]

\medskip
{\small\bfseries [Concept 2]:}\; [one-line formula/definition]

\medskip
{\small\bfseries [Technique/recipe, if applicable]:}
\step{\ding{182}}{[step]}
\step{\ding{183}}{[step]}

\begin{keybox}
\centering\small\bfseries\color{redline} Must remember: [the 2-4
things that, if forgotten, would sink an otherwise-good answer.]
\end{keybox}
\end{rc}

% ============================================================
% ASSESSMENT-MAPPING TABLE — every real item touching this chapter.
% This is what tells the learner exactly what to go practice next.
% ============================================================
\topicsection{[X.N+1] Assessment Mapping}

\begin{rc}
{\small
\renewcommand{\arraystretch}{1.3}
\begin{tabular}{@{}lp{10.5cm}@{}}
\toprule
\textbf{Source} & \textbf{Topic / key technique tested} \\
\midrule
[assessment:S004, 2021 exam Q3] & [what it tested, in one line] \\
[source:S002, Worksheet 4 Task 2] & [what it tested] \\
[assessment:S005, Assignment 2 Q1] & [what it tested] \\
\bottomrule
\end{tabular}}

\begin{keybox}
\centering\small\bfseries\color{redline} [One-sentence summary of the
highest-probability item type for the next assessment, based on
recurrence across the sources above.]
\end{keybox}
\end{rc}

\end{document}
